\documentclass{amsart}
% For dealing with references we use the comment environment
\usepackage{verbatim}
\newenvironment{reference}{\comment}{\endcomment}
%\newenvironment{reference}{}{}
\newenvironment{slogan}{\comment}{\endcomment}
\newenvironment{history}{\comment}{\endcomment}

% For commutative diagrams we use Xy-pic
\usepackage[all]{xy}

% We use 2cell for 2-commutative diagrams.
\xyoption{2cell}
\UseAllTwocells

% We use multicol for the list of chapters between chapters
\usepackage{multicol}

% This is generally recommended for better output
\usepackage{lmodern}
\usepackage[T1]{fontenc}

% For cross-file-references

% Package for hypertext links:
\usepackage{hyperref}

% For any local file, say "hello.tex" you want to link to please

% Theorem environments.
%
\theoremstyle{plain}
\newtheorem{theorem}[subsection]{Theorem}
\newtheorem{proposition}[subsection]{Proposition}
\newtheorem{lemma}[subsection]{Lemma}

\theoremstyle{definition}
\newtheorem{definition}[subsection]{Definition}
\newtheorem{example}[subsection]{Example}
\newtheorem{exercise}[subsection]{Exercise}
\newtheorem{situation}[subsection]{Situation}

\theoremstyle{remark}
\newtheorem{remark}[subsection]{Remark}
\newtheorem{remarks}[subsection]{Remarks}

\numberwithin{equation}{subsection}

% Macros
%
\def\lim{\mathop{\mathrm{lim}}\nolimits}
\def\colim{\mathop{\mathrm{colim}}\nolimits}
\def\Spec{\mathop{\mathrm{Spec}}}
\def\Hom{\mathop{\mathrm{Hom}}\nolimits}
\def\Ext{\mathop{\mathrm{Ext}}\nolimits}
\def\SheafHom{\mathop{\mathcal{H}\!\mathit{om}}\nolimits}
\def\SheafExt{\mathop{\mathcal{E}\!\mathit{xt}}\nolimits}
\def\Sch{\mathit{Sch}}
\def\Mor{\mathop{\mathrm{Mor}}\nolimits}
\def\Ob{\mathop{\mathrm{Ob}}\nolimits}
\def\Sh{\mathop{\mathit{Sh}}\nolimits}
\def\NL{\mathop{N\!L}\nolimits}
\def\CH{\mathop{\mathrm{CH}}\nolimits}
\def\proetale{{pro\text{-}\acute{e}tale}}
\def\etale{{\acute{e}tale}}
\def\QCoh{\mathit{QCoh}}
\def\Ker{\mathop{\mathrm{Ker}}}
\def\Im{\mathop{\mathrm{Im}}}
\def\Coker{\mathop{\mathrm{Coker}}}
\def\Coim{\mathop{\mathrm{Coim}}}

% Boxtimes
%
\DeclareMathSymbol{\boxtimes}{\mathbin}{AMSa}{"02}

%
% Macros for moduli stacks/spaces
%
\def\QCohstack{\mathcal{QC}\!\mathit{oh}}
\def\Cohstack{\mathcal{C}\!\mathit{oh}}
\def\Spacesstack{\mathcal{S}\!\mathit{paces}}
\def\Quotfunctor{\mathrm{Quot}}
\def\Hilbfunctor{\mathrm{Hilb}}
\def\Curvesstack{\mathcal{C}\!\mathit{urves}}
\def\Polarizedstack{\mathcal{P}\!\mathit{olarized}}
\def\Complexesstack{\mathcal{C}\!\mathit{omplexes}}
% \Pic is the operator that assigns to X its picard group, usage \Pic(X)
% \Picardstack_{X/B} denotes the Picard stack of X over B
% \Picardfunctor_{X/B} denotes the Picard functor of X over B
\def\Pic{\mathop{\mathrm{Pic}}\nolimits}
\def\Picardstack{\mathcal{P}\!\mathit{ic}}
\def\Picardfunctor{\mathrm{Pic}}
\def\Deformationcategory{\mathcal{D}\!\mathit{ef}}

\setlength{\emergencystretch}{3em}
\begin{document}
\title[Stacks Project, Tag 0EC3]{314 --- Injectivity of residue-supported Frobenius-stable modules over regular local rings}
\maketitle
\noindent Copyright (C) 2005--2025 Johan de Jong. Stacks Project authors. Source: \href{https://stacks.math.columbia.edu/tag/0EC3}{Tag 0EC3}.

\noindent Editorial difficulty note (not author text). Difficulty H2: The proof computes all Frobenius-power torsion layers using flatness of iterated Frobenius and shows these layers are free over truncated parameter quotients. It obtains parameter divisibility, then verifies Frobenius stability passes to parameter torsion and inducts on dimension. This is a distinct positive-characteristic classification task, not a parameter change of the derivation proof..

\noindent Editorial provenance note (not author text). History: excerpt from revision \texttt{a0444\allowbreak 6e57e\allowbreak c1fbc\allowbreak 252a8\allowbreak 71afc\allowbreak ec775\allowbreak 2fb28\allowbreak 07b14}, local-cohomology.tex, lines 3857--3951. Reference presentation was adapted; original-author context and core support blocks were retained or added. The mathematical wording is unchanged. Licensed under GNU FDL 1.2 or later; no invariant sections or cover texts. See ../licenses/COPYING. Context blocks reproduce author definitions and supporting results; they are not additional counted problems.

\begin{lemma}
\label{lemma-structure-torsion-Frobenius-regular}
\begin{reference}
Follows from \href{https://stacks.math.columbia.edu/bibliography}{Bibliography: Huneke-Sharp (Corollary 3.6)} with a
little bit of work. Also follows directly from
\href{https://stacks.math.columbia.edu/bibliography}{Bibliography: Lyubeznik2 (Theorem 1.4)}.
\end{reference}
Let $p$ be a prime number. Let $(A, \mathfrak m, k)$
be a regular local ring with $p = 0$. Denote $F : A \to A$, $a \mapsto a^p$
be the Frobenius endomorphism. Let $M$ be a $\mathfrak m$-power torsion module
such that $M \otimes_{A, F} A \cong M$. Then $M$ is isomorphic to a direct sum
of copies of the injective hull $E$ of $k$.
\end{lemma}

\begin{proof}
Choose a set $J$ and an $A$-module homorphism
$\varphi : M \to \bigoplus_{j \in J} E$ which maps
$M[\mathfrak m]$ isomorphically onto
$(\bigoplus_{j \in J} E)[\mathfrak m] = \bigoplus_{j \in J} k$.
We claim that $\varphi$ is an isomorphism, i.e., bijective.

\medskip\noindent
Injective. Let $z \in M$ be nonzero. Since $M$ is $\mathfrak m$-power torsion
we can choose an element $f \in A$ such that $fz \in M[\mathfrak m]$ and
$fz \not = 0$. Then $\varphi(fz) = f\varphi(z)$ is nonzero, hence
$\varphi(z)$ is nonzero.

\medskip\noindent
Surjective. Recall that $F$ is flat, see
Lemma \href{https://stacks.math.columbia.edu/tag/0EC0}{lemma frobenius flat regular (Tag 0EC0)}.
Let $x_1, \ldots, x_d$ be a minimal system of generators of
$\mathfrak m$. Denote
$$
M_n = M[x_1^{p^n}, \ldots, x_d^{p^n}]
$$
the submodule of $M$ consisting of elements killed by
$x_1^{p^n}, \ldots, x_d^{p^n}$. So $M_0 = M[\mathfrak m]$
is a vector space over $k$. Also $M = \bigcup M_n$ by our
assumption that $M$ is $\mathfrak m$-power torsion. Since $F^n$ is flat and
$F^n(x_i) = x_i^{p^n}$ we have
$$
M_n \cong (M \otimes_{A, F^n} A)[x_1^{p^n}, \ldots, x_d^{p^n}] =
M[x_1, \ldots, x_d] \otimes_{A, F^n} A =
M_0 \otimes_k A/(x_1^{p^n}, \ldots, x_d^{p^n})
$$
Thus $M_n$ is free over $A/(x_1^{p^n}, \ldots, x_d^{p^n})$.
A computation shows that every element of $A/(x_1^{p^n}, \ldots, x_d^{p^n})$
annihilated by $x_1^{p^n - 1}$ is divisible by $x_1$; for example
you can use that $A/(x_1^{p^n}, \ldots, x_d^{p^n}) \cong
k[x_1, \ldots, x_d]/(x_1^{p^n}, \ldots, x_d^{p^n})$ by Algebra, Lemma
\href{https://stacks.math.columbia.edu/tag/0C0S}{lemma regular complete containing coefficient field (Tag 0C0S)}.
Thus the same is true for every element of $M_n$.
Since every element of $M$ is in $M_n$ for all $n \gg 0$
and since every element of $M$ is killed by some power of
$x_1$, we conclude that $M$ is $x_1$-divisible.

\medskip\noindent
Let $x = x_1$. Above we have seen that $M$ is $x$-divisible.
If $\varphi$ is not surjective, then we can choose
$e \in \bigoplus_{j \in J} E$ not in $M$.
Arguing as above we may assume $\mathfrak m e \subset M$,
in particular $x e \in M$. There exists an element
$z_1 \in M$ with $x z_1 = x e$. Hence
$x(z_1 - e) = 0$. Replacing $e$ by $e - z_1$
we may assume $e$ is annihilated by $x$.
Thus it suffices to prove that
$$
\varphi[x] :
M[x]
\longrightarrow
\left(\bigoplus\nolimits_{j \in J} E\right)[x] =
\bigoplus\nolimits_{j \in J} E[x]
$$
is surjective. If $d = 1$, this is true by construction of $\varphi$.
If $d > 1$, then we observe that $E[x]$ is the injective hull
of the residue field of the regular ring $A/xA$, see
Dualizing Complexes, Lemma \href{https://stacks.math.columbia.edu/tag/08Z2}{lemma quotient (Tag 08Z2)}.
Observe that $M[x]$ as a module over $A/xA$
is $\mathfrak m/(x)$-power torsion and we have
\begin{align*}
M[x] \otimes_{A/xA, F} A/xA
& = M[x] \otimes_{A, F} A \otimes_A A/xA \\
& = (M \otimes_{A, F} A)[x^p] \otimes_A A/xA \\
& \cong M[x^p] \otimes_A A/xA
\end{align*}
Argue using flatness of $F$ as before.
We claim that $M[x^p] \otimes_A A/xA \to M[x]$,
$z \otimes 1 \mapsto x^{p - 1}z$ is an isomorphism.
This can be seen by proving it for
each of the modules $M_n$, $n > 0$ defined above
where it follows by the same result for
$A/(x_1^{p^n}, \ldots, x_d^{p^n})$ and $x = x_1$.
Thus by induction on $\dim(A)$ we conclude that $\varphi[x]$
is surjective as desired.
\end{proof}
\end{document}
